% Reproducibility Checklist
\section{Reproducibility Checklist}

\begin{enumerate}
    \item This paper:
    \begin{itemize}
        \item Includes a conceptual outline and/or pseudocode description of AI methods introduced (\textbf{\underline{yes}}/partial/no/NA)
        \item Clearly delineates statements that are opinions, hypotheses, and speculation from objective facts and results (\highlight{yes}/no)
        \item Provides well-marked pedagogical references for less-familiar readers to gain background necessary to replicate the paper (\highlight{yes}/no)
    \end{itemize}

    \item Does this paper make theoretical contributions? (yes/\highlight{no}) \\


    \item Does this paper rely on one or more datasets? (\highlight{yes}/no) \\
    If yes, please complete the list below:
    \begin{itemize}
        \item A motivation is given for why the experiments are conducted on the selected datasets. (\highlight{yes}/partial/no/NA)
        \item All novel datasets introduced in this paper are included in a data appendix. (yes/partial/\highlight{no}/NA)
        \item All novel datasets introduced in this paper will be made publicly available upon publication with a license allowing free research use. (\highlight{yes}/partial/no/{NA})
        \item All datasets drawn from the existing literature are accompanied by appropriate citations. (\highlight{yes}/no/NA)
        \item All datasets drawn from the existing literature are publicly available. (\highlight{yes}/partial/no/NA)
        \item Datasets that are not publicly available are described in detail, with justification. (\highlight{yes}/partial/no/NA)
    \end{itemize}

    \item Does this paper include computational experiments? (\highlight{yes}/no) \\
    If yes, please complete the list below:
    \begin{itemize}
        \item This paper states the number and range of values tried per (hyper-) parameter during development of the paper, along with the criterion used for selecting the final parameter setting. (\highlight{yes}/partial/no/NA)
        \item Any code required for pre-processing data is included in the appendix. (yes/partial/\highlight{no})
        \item All source code required for conducting and analyzing the experiments is included in a code appendix. (\highlight{yes}/partial/no)
        \item All source code required for conducting and analyzing the experiments will be made publicly available upon publication of the paper with a license that allows free usage for research purposes. (\highlight{yes}/partial/no)
        \item All source code implementing new methods have comments detailing the implementation, with references to the paper where each step comes from (\highlight{yes}/partial/no)
        \item If an algorithm depends on randomness, then the method used for setting seeds is described in a way sufficient to allow replication of results. (yes/partial/no/\highlight{NA})
        \item This paper specifies the computing infrastructure used for running experiments (hardware and software), including GPU/CPU models; amount of memory; operating system; names and versions of relevant software libraries and frameworks. (\highlight{yes}/partial/no)
        \item This paper states the number of algorithm runs used to compute each reported result. (\highlight{yes}/partial/no)
        \item Analysis of experiments goes beyond single-dimensional summaries of performance (e.g., average; median) to include measures of variation, confidence, or other distributional information. (yes/\highlight{no})
        \item The significance of any improvement or decrease in performance is judged using appropriate statistical tests (e.g., Wilcoxon signed-rank). (yes/\highlight{no})
        \item Significance of performance differences is assessed with statistical tests. (yes/partial/\highlight{no})
        \item This paper lists all final (hyper-)parameters used for each model/algorithm in the paper’s experiments. (yes/\highlight{partial}/no/NA)
    \end{itemize}
\end{enumerate}



